%%%%%%%%%%%%%%%%%%%%%%%%%%%%%%%%%%%%%%%%%%%%%%%%%%%%%%%%%%%%
% 13.7 NETWORK SCIENCE
% The Composition of Advanced Mathematics
%%%%%%%%%%%%%%%%%%%%%%%%%%%%%%%%%%%%%%%%%%%%%%%%%%%%%%%%%%%%

\section{Network Science}
\label{sec:network-science}

\prerequisite{
Graph theory, linear algebra, probability, differential equations,
dynamical systems, optimization, statistics, and mathematical modeling.
}

\learningobjectives{
By the end of this section, the reader should be able to:
\begin{itemize}
    \item explain the purpose of network science;
    \item distinguish graphs from physical networks;
    \item define vertices, edges, directed edges, and weighted edges;
    \item construct adjacency matrices;
    \item define degree, in-degree, and out-degree;
    \item understand degree distributions;
    \item recognize sparse and dense networks;
    \item define walks, trails, paths, and cycles;
    \item understand connected components;
    \item distinguish weak and strong connectivity;
    \item define graph distance;
    \item understand diameter and average path length;
    \item understand clustering coefficients;
    \item recognize triangles and transitivity;
    \item understand local and global network structure;
    \item define centrality;
    \item compute degree centrality;
    \item understand closeness centrality;
    \item understand betweenness centrality;
    \item understand eigenvector centrality;
    \item recognize PageRank conceptually;
    \item understand Katz centrality conceptually;
    \item understand hubs and authorities conceptually;
    \item define the graph Laplacian;
    \item understand Laplacian quadratic forms;
    \item recognize spectral graph theory;
    \item understand algebraic connectivity;
    \item recognize Fiedler vectors conceptually;
    \item understand spectral clustering conceptually;
    \item define random graphs;
    \item understand Erdős--Rényi networks;
    \item recognize expected degree;
    \item understand degree distributions in random graphs;
    \item recognize giant-component transitions conceptually;
    \item understand small-world networks;
    \item recognize Watts--Strogatz models;
    \item understand scale-free networks;
    \item recognize preferential attachment;
    \item understand heavy-tailed degree distributions;
    \item distinguish degree distributions from power laws;
    \item understand configuration models;
    \item recognize assortative and disassortative mixing;
    \item understand homophily conceptually;
    \item understand community structure;
    \item define modularity conceptually;
    \item understand community detection;
    \item recognize hierarchical networks;
    \item understand bipartite networks;
    \item recognize projections of bipartite networks;
    \item understand multilayer and multiplex networks;
    \item recognize temporal networks;
    \item understand network robustness;
    \item recognize node and edge failures;
    \item understand percolation conceptually;
    \item recognize cascading failures;
    \item understand epidemic spreading on networks;
    \item recognize threshold effects;
    \item understand bond and site percolation conceptually;
    \item understand diffusion on networks;
    \item formulate Laplacian diffusion equations;
    \item understand random walks on graphs;
    \item construct transition matrices;
    \item understand stationary distributions;
    \item recognize Markov-chain connections;
    \item understand consensus dynamics;
    \item recognize synchronization conceptually;
    \item understand coupled dynamical systems on networks;
    \item recognize network controllability conceptually;
    \item understand flow networks;
    \item recognize shortest paths and routing;
    \item understand maximum flow and minimum cut conceptually;
    \item recognize transportation and communication networks;
    \item understand social-network modeling;
    \item recognize biological networks;
    \item understand ecological networks;
    \item recognize financial networks;
    \item understand infrastructure networks;
    \item recognize dependency networks;
    \item understand network inference from data;
    \item recognize missing and uncertain edges;
    \item understand link prediction conceptually;
    \item recognize network sampling bias;
    \item understand null models;
    \item recognize network motifs;
    \item understand motif significance conceptually;
    \item recognize network entropy conceptually;
    \item understand dynamical processes on networks;
    \item connect network science with graph theory, probability,
          dynamical systems, control, statistics, optimization,
          biology, finance, and economics.
\end{itemize}
}

%%%%%%%%%%%%%%%%%%%%%%%%%%%%%%%%%%%%%%%%%%%%%%%%%%%%%%%%%%%%
\subsection{What Is Network Science?}
\label{subsec:what-is-network-science}
%%%%%%%%%%%%%%%%%%%%%%%%%%%%%%%%%%%%%%%%%%%%%%%%%%%%%%%%%%%%

Network science studies systems whose behavior depends strongly on
relationships among interacting components.

A network is commonly represented by a graph

\[
\boxed{
G=(V,E),
}
\]

where

\[
V
=
\text{set of vertices or nodes}
\]

and

\[
E
=
\text{set of edges or links}.
\]

The graph provides the mathematical structure.

The network gives the graph an applied interpretation.

%%%%%%%%%%%%%%%%%%%%%%%%%%%%%%%%%%%%%%%%%%%%%%%%%%%%%%%%%%%%
\subsection{Examples of Networks}
\label{subsec:network-examples}
%%%%%%%%%%%%%%%%%%%%%%%%%%%%%%%%%%%%%%%%%%%%%%%%%%%%%%%%%%%%

Examples include:

\[
\boxed{
\text{social networks},
}
\]

\[
\boxed{
\text{transportation networks},
}
\]

\[
\boxed{
\text{communication networks},
}
\]

\[
\boxed{
\text{power grids},
}
\]

\[
\boxed{
\text{biological networks},
}
\]

and

\[
\boxed{
\text{financial networks}.
}
\]

The same graph-theoretic tools can describe very different applications.

%%%%%%%%%%%%%%%%%%%%%%%%%%%%%%%%%%%%%%%%%%%%%%%%%%%%%%%%%%%%
\subsection{Undirected Networks}
\label{subsec:undirected-networks}
%%%%%%%%%%%%%%%%%%%%%%%%%%%%%%%%%%%%%%%%%%%%%%%%%%%%%%%%%%%%

In an undirected network, an edge

\[
\{i,j\}
\]

has no direction.

Examples include symmetric friendship relationships and some physical
connections.

The adjacency matrix satisfies

\[
\boxed{
A_{ij}=A_{ji}.
}
\]

Thus \(A\) is symmetric.

%%%%%%%%%%%%%%%%%%%%%%%%%%%%%%%%%%%%%%%%%%%%%%%%%%%%%%%%%%%%
\subsection{Directed Networks}
\label{subsec:directed-networks}
%%%%%%%%%%%%%%%%%%%%%%%%%%%%%%%%%%%%%%%%%%%%%%%%%%%%%%%%%%%%

A directed network contains ordered edges

\[
\boxed{
(i,j).
}
\]

The direction may represent:

\[
\boxed{
\text{information flow},
}
\]

\[
\boxed{
\text{citation},
}
\]

\[
\boxed{
\text{transport},
}
\]

or

\[
\boxed{
\text{influence}.
}
\]

In general,

\[
A_{ij}
\neq
A_{ji}.
\]

%%%%%%%%%%%%%%%%%%%%%%%%%%%%%%%%%%%%%%%%%%%%%%%%%%%%%%%%%%%%
\subsection{Weighted Networks}
\label{subsec:weighted-networks}
%%%%%%%%%%%%%%%%%%%%%%%%%%%%%%%%%%%%%%%%%%%%%%%%%%%%%%%%%%%%

Edges may carry weights

\[
\boxed{
w_{ij}.
}
\]

A weight can represent:

\[
\boxed{
\text{distance},
}
\]

\[
\boxed{
\text{capacity},
}
\]

\[
\boxed{
\text{frequency},
}
\]

\[
\boxed{
\text{cost},
}
\]

or

\[
\boxed{
\text{interaction strength}.
}
\]

The meaning of larger weight depends on the application.

%%%%%%%%%%%%%%%%%%%%%%%%%%%%%%%%%%%%%%%%%%%%%%%%%%%%%%%%%%%%
\subsection{Adjacency Matrix}
\label{subsec:network-adjacency-matrix}
%%%%%%%%%%%%%%%%%%%%%%%%%%%%%%%%%%%%%%%%%%%%%%%%%%%%%%%%%%%%

For a simple unweighted graph with \(n\) nodes, define

\[
\boxed{
A_{ij}
=
\begin{cases}
1, & \text{if nodes }i\text{ and }j\text{ are adjacent},\\
0, & \text{otherwise}.
\end{cases}
}
\]

For a simple undirected graph without self-loops,

\[
A_{ii}=0
\]

and

\[
A=A^T.
\]

%%%%%%%%%%%%%%%%%%%%%%%%%%%%%%%%%%%%%%%%%%%%%%%%%%%%%%%%%%%%
\subsection{Worked Example: Adjacency Matrix}
\label{subsec:worked-network-adjacency}
%%%%%%%%%%%%%%%%%%%%%%%%%%%%%%%%%%%%%%%%%%%%%%%%%%%%%%%%%%%%

\begin{workedexamplebox}{A Three-Node Path}

Consider

\[
1
\longleftrightarrow
2
\longleftrightarrow
3.
\]

The adjacency matrix is

\[
A
=
\begin{bmatrix}
0 & 1 & 0\\
1 & 0 & 1\\
0 & 1 & 0
\end{bmatrix}.
\]

\begin{answerbox}
\[
\boxed{
A
=
\begin{bmatrix}
0 & 1 & 0\\
1 & 0 & 1\\
0 & 1 & 0
\end{bmatrix}.
}
\]
\end{answerbox}

\end{workedexamplebox}

%%%%%%%%%%%%%%%%%%%%%%%%%%%%%%%%%%%%%%%%%%%%%%%%%%%%%%%%%%%%
\subsection{Degree}
\label{subsec:network-degree}
%%%%%%%%%%%%%%%%%%%%%%%%%%%%%%%%%%%%%%%%%%%%%%%%%%%%%%%%%%%%

For an undirected graph, degree of node \(i\) is

\[
\boxed{
k_i
=
\sum_jA_{ij}.
}
\]

It counts the number of neighbors of \(i\).

The average degree is

\[
\boxed{
\langle k\rangle
=
\frac1n
\sum_{i=1}^{n}k_i.
}
\]

%%%%%%%%%%%%%%%%%%%%%%%%%%%%%%%%%%%%%%%%%%%%%%%%%%%%%%%%%%%%
\subsection{Handshaking Relation}
\label{subsec:network-handshaking}
%%%%%%%%%%%%%%%%%%%%%%%%%%%%%%%%%%%%%%%%%%%%%%%%%%%%%%%%%%%%

For an undirected graph,

\[
\boxed{
\sum_{i=1}^{n}k_i
=
2|E|.
}
\]

Every edge contributes one unit of degree to each endpoint.

Therefore,

\[
\boxed{
\langle k\rangle
=
\frac{2|E|}{|V|}.
}
\]

%%%%%%%%%%%%%%%%%%%%%%%%%%%%%%%%%%%%%%%%%%%%%%%%%%%%%%%%%%%%
\subsection{In-Degree and Out-Degree}
\label{subsec:network-in-out-degree}
%%%%%%%%%%%%%%%%%%%%%%%%%%%%%%%%%%%%%%%%%%%%%%%%%%%%%%%%%%%%

For a directed graph, define out-degree

\[
\boxed{
k_i^{\mathrm{out}}
=
\sum_jA_{ij},
}
\]

and in-degree

\[
\boxed{
k_i^{\mathrm{in}}
=
\sum_jA_{ji}.
}
\]

A node may have very different incoming and outgoing connectivity.

%%%%%%%%%%%%%%%%%%%%%%%%%%%%%%%%%%%%%%%%%%%%%%%%%%%%%%%%%%%%
\subsection{Strength}
\label{subsec:network-strength}
%%%%%%%%%%%%%%%%%%%%%%%%%%%%%%%%%%%%%%%%%%%%%%%%%%%%%%%%%%%%

For a weighted network, weighted degree is often called strength:

\[
\boxed{
s_i
=
\sum_jw_{ij}.
}
\]

Degree measures the number of connections.

Strength measures their total weight.

%%%%%%%%%%%%%%%%%%%%%%%%%%%%%%%%%%%%%%%%%%%%%%%%%%%%%%%%%%%%
\subsection{Degree Distribution}
\label{subsec:degree-distribution}
%%%%%%%%%%%%%%%%%%%%%%%%%%%%%%%%%%%%%%%%%%%%%%%%%%%%%%%%%%%%

The degree distribution is

\[
\boxed{
P(k)
=
\Pr(K=k),
}
\]

or empirically,

\[
\boxed{
P(k)
=
\frac{
\text{number of nodes with degree }k
}{
n
}.
}
\]

It summarizes variation in local connectivity.

%%%%%%%%%%%%%%%%%%%%%%%%%%%%%%%%%%%%%%%%%%%%%%%%%%%%%%%%%%%%
\subsection{Sparse Networks}
\label{subsec:sparse-networks}
%%%%%%%%%%%%%%%%%%%%%%%%%%%%%%%%%%%%%%%%%%%%%%%%%%%%%%%%%%%%

Many real networks are sparse.

For an undirected simple network, the maximum number of edges is

\[
\binom{n}{2}.
\]

A network is sparse when its number of edges is much smaller than this
maximum.

Frequently,

\[
\boxed{
|E|=O(n)
}
\]

rather than \(O(n^2)\).

%%%%%%%%%%%%%%%%%%%%%%%%%%%%%%%%%%%%%%%%%%%%%%%%%%%%%%%%%%%%
\subsection{Walks and Paths}
\label{subsec:network-walks-paths}
%%%%%%%%%%%%%%%%%%%%%%%%%%%%%%%%%%%%%%%%%%%%%%%%%%%%%%%%%%%%

A walk is a sequence of adjacent vertices.

A path typically means a walk without repeated vertices.

For example,

\[
\boxed{
1\to3\to5\to7
}
\]

is a path if the corresponding edges exist and no vertex repeats.

Paths determine reachability and distance.

%%%%%%%%%%%%%%%%%%%%%%%%%%%%%%%%%%%%%%%%%%%%%%%%%%%%%%%%%%%%
\subsection{Graph Distance}
\label{subsec:graph-distance-network}
%%%%%%%%%%%%%%%%%%%%%%%%%%%%%%%%%%%%%%%%%%%%%%%%%%%%%%%%%%%%

The distance between vertices \(i\) and \(j\) is the length of a shortest
path connecting them:

\[
\boxed{
d(i,j).
}
\]

If no path exists, one may treat the distance as infinite.

%%%%%%%%%%%%%%%%%%%%%%%%%%%%%%%%%%%%%%%%%%%%%%%%%%%%%%%%%%%%
\subsection{Connected Networks}
\label{subsec:connected-networks}
%%%%%%%%%%%%%%%%%%%%%%%%%%%%%%%%%%%%%%%%%%%%%%%%%%%%%%%%%%%%

An undirected graph is connected if every pair of vertices is linked by a
path.

Otherwise it decomposes into connected components.

A connected component is a maximal connected subgraph.

%%%%%%%%%%%%%%%%%%%%%%%%%%%%%%%%%%%%%%%%%%%%%%%%%%%%%%%%%%%%
\subsection{Strong and Weak Connectivity}
\label{subsec:strong-weak-connectivity}
%%%%%%%%%%%%%%%%%%%%%%%%%%%%%%%%%%%%%%%%%%%%%%%%%%%%%%%%%%%%

For directed graphs, strong connectivity means every node can reach every
other node by directed paths.

Weak connectivity means the underlying undirected graph is connected after
edge directions are ignored.

%%%%%%%%%%%%%%%%%%%%%%%%%%%%%%%%%%%%%%%%%%%%%%%%%%%%%%%%%%%%
\subsection{Eccentricity}
\label{subsec:network-eccentricity}
%%%%%%%%%%%%%%%%%%%%%%%%%%%%%%%%%%%%%%%%%%%%%%%%%%%%%%%%%%%%

The eccentricity of node \(i\) is

\[
\boxed{
e(i)
=
\max_jd(i,j)
}
\]

within a connected graph.

It measures the distance from \(i\) to its farthest node.

%%%%%%%%%%%%%%%%%%%%%%%%%%%%%%%%%%%%%%%%%%%%%%%%%%%%%%%%%%%%
\subsection{Diameter}
\label{subsec:network-diameter}
%%%%%%%%%%%%%%%%%%%%%%%%%%%%%%%%%%%%%%%%%%%%%%%%%%%%%%%%%%%%

The diameter is

\[
\boxed{
D
=
\max_{i,j}
d(i,j).
}
\]

It is the greatest shortest-path distance between connected nodes.

%%%%%%%%%%%%%%%%%%%%%%%%%%%%%%%%%%%%%%%%%%%%%%%%%%%%%%%%%%%%
\subsection{Average Path Length}
\label{subsec:average-path-length}
%%%%%%%%%%%%%%%%%%%%%%%%%%%%%%%%%%%%%%%%%%%%%%%%%%%%%%%%%%%%

For a connected undirected network, average shortest-path length can be
written

\[
\boxed{
\ell
=
\frac{1}{n(n-1)}
\sum_{i\neq j}
d(i,j).
}
\]

Small average path length indicates that nodes are typically only a few steps
apart.

%%%%%%%%%%%%%%%%%%%%%%%%%%%%%%%%%%%%%%%%%%%%%%%%%%%%%%%%%%%%
\subsection{Triangles}
\label{subsec:network-triangles}
%%%%%%%%%%%%%%%%%%%%%%%%%%%%%%%%%%%%%%%%%%%%%%%%%%%%%%%%%%%%

A triangle occurs when three nodes are pairwise connected.

Triangles measure local closure.

In a simple undirected graph, the number of triangles is

\[
\boxed{
\frac{1}{6}
\operatorname{tr}(A^3).
}
\]

Each triangle is counted six times in \(\operatorname{tr}(A^3)\).

%%%%%%%%%%%%%%%%%%%%%%%%%%%%%%%%%%%%%%%%%%%%%%%%%%%%%%%%%%%%
\subsection{Local Clustering Coefficient}
\label{subsec:local-clustering}
%%%%%%%%%%%%%%%%%%%%%%%%%%%%%%%%%%%%%%%%%%%%%%%%%%%%%%%%%%%%

Suppose node \(i\) has degree \(k_i\geq2\).

If \(e_i\) edges exist among its neighbors, define

\[
\boxed{
C_i
=
\frac{2e_i}
{k_i(k_i-1)}.
}
\]

This measures the fraction of possible neighbor-neighbor connections that
actually occur.

%%%%%%%%%%%%%%%%%%%%%%%%%%%%%%%%%%%%%%%%%%%%%%%%%%%%%%%%%%%%
\subsection{Worked Example: Clustering}
\label{subsec:worked-network-clustering}
%%%%%%%%%%%%%%%%%%%%%%%%%%%%%%%%%%%%%%%%%%%%%%%%%%%%%%%%%%%%

\begin{workedexamplebox}{Clustering of a Node}

Suppose a node has

\[
k_i=4
\]

neighbors.

There are

\[
\binom42=6
\]

possible edges among those neighbors.

If \(3\) such edges exist, then

\[
C_i
=
\frac{3}{6}.
\]

Therefore,

\begin{answerbox}
\[
\boxed{
C_i=\frac12.
}
\]
\end{answerbox}

\end{workedexamplebox}

%%%%%%%%%%%%%%%%%%%%%%%%%%%%%%%%%%%%%%%%%%%%%%%%%%%%%%%%%%%%
\subsection{Global Transitivity}
\label{subsec:network-transitivity}
%%%%%%%%%%%%%%%%%%%%%%%%%%%%%%%%%%%%%%%%%%%%%%%%%%%%%%%%%%%%

A common global clustering measure is

\[
\boxed{
C
=
\frac{
3\times\text{number of triangles}
}{
\text{number of connected triples}
}.
}
\]

This measures how frequently two-step relationships close into triangles.

%%%%%%%%%%%%%%%%%%%%%%%%%%%%%%%%%%%%%%%%%%%%%%%%%%%%%%%%%%%%
\subsection{Centrality}
\label{subsec:network-centrality}
%%%%%%%%%%%%%%%%%%%%%%%%%%%%%%%%%%%%%%%%%%%%%%%%%%%%%%%%%%%%

Centrality measures structural importance of a node.

Different centralities answer different questions.

A node may be important because it has:

\[
\boxed{
\text{many neighbors},
}
\]

\[
\boxed{
\text{short distances to others},
}
\]

\[
\boxed{
\text{control over shortest paths},
}
\]

or

\[
\boxed{
\text{connections to other important nodes}.
}
\]

%%%%%%%%%%%%%%%%%%%%%%%%%%%%%%%%%%%%%%%%%%%%%%%%%%%%%%%%%%%%
\subsection{Degree Centrality}
\label{subsec:degree-centrality}
%%%%%%%%%%%%%%%%%%%%%%%%%%%%%%%%%%%%%%%%%%%%%%%%%%%%%%%%%%%%

For an undirected simple network,

\[
\boxed{
C_D(i)
=
\frac{k_i}{n-1}.
}
\]

This normalizes degree to the interval

\[
0\leq C_D(i)\leq1.
\]

Degree centrality measures immediate local connectivity.

%%%%%%%%%%%%%%%%%%%%%%%%%%%%%%%%%%%%%%%%%%%%%%%%%%%%%%%%%%%%
\subsection{Closeness Centrality}
\label{subsec:closeness-centrality}
%%%%%%%%%%%%%%%%%%%%%%%%%%%%%%%%%%%%%%%%%%%%%%%%%%%%%%%%%%%%

For a connected graph, closeness centrality is often

\[
\boxed{
C_C(i)
=
\frac{n-1}
{\sum_{j\neq i}d(i,j)}.
}
\]

A node has high closeness when its average shortest-path distance to other
nodes is small.

%%%%%%%%%%%%%%%%%%%%%%%%%%%%%%%%%%%%%%%%%%%%%%%%%%%%%%%%%%%%
\subsection{Betweenness Centrality}
\label{subsec:betweenness-centrality}
%%%%%%%%%%%%%%%%%%%%%%%%%%%%%%%%%%%%%%%%%%%%%%%%%%%%%%%%%%%%

Let

\[
\sigma_{st}
\]

be the number of shortest paths from \(s\) to \(t\), and

\[
\sigma_{st}(i)
\]

the number passing through node \(i\).

Then

\[
\boxed{
C_B(i)
=
\sum_{\substack{s\neq i\neq t\\s\neq t}}
\frac{\sigma_{st}(i)}
{\sigma_{st}}.
}
\]

Normalization conventions vary.

%%%%%%%%%%%%%%%%%%%%%%%%%%%%%%%%%%%%%%%%%%%%%%%%%%%%%%%%%%%%
\subsection{Interpretation of Betweenness}
\label{subsec:betweenness-interpretation}
%%%%%%%%%%%%%%%%%%%%%%%%%%%%%%%%%%%%%%%%%%%%%%%%%%%%%%%%%%%%

High-betweenness nodes can act as:

\[
\boxed{
\text{bridges},
}
\]

\[
\boxed{
\text{brokers},
}
\]

\[
\boxed{
\text{transport bottlenecks},
}
\]

or

\[
\boxed{
\text{communication intermediaries}.
}
\]

A node can have modest degree but high betweenness.

%%%%%%%%%%%%%%%%%%%%%%%%%%%%%%%%%%%%%%%%%%%%%%%%%%%%%%%%%%%%
\subsection{Eigenvector Centrality}
\label{subsec:eigenvector-centrality}
%%%%%%%%%%%%%%%%%%%%%%%%%%%%%%%%%%%%%%%%%%%%%%%%%%%%%%%%%%%%

Eigenvector centrality assigns higher importance to nodes connected to other
important nodes.

Let \(\mathbf{x}\) satisfy

\[
\boxed{
A\mathbf{x}
=
\lambda\mathbf{x}.
}
\]

The centrality vector is typically associated with the dominant eigenvalue
under appropriate nonnegative-network conditions.

Thus

\[
\boxed{
x_i
\propto
\sum_jA_{ij}x_j.
}
\]

%%%%%%%%%%%%%%%%%%%%%%%%%%%%%%%%%%%%%%%%%%%%%%%%%%%%%%%%%%%%
\subsection{PageRank}
\label{subsec:pagerank}
%%%%%%%%%%%%%%%%%%%%%%%%%%%%%%%%%%%%%%%%%%%%%%%%%%%%%%%%%%%%

PageRank modifies eigenvector-style centrality for directed networks by
including random transitions.

A simplified form is

\[
\boxed{
\mathbf{p}
=
\alpha P^T\mathbf{p}
+
(1-\alpha)\mathbf{v},
}
\]

where:

\[
P
=
\text{transition matrix},
\]

\[
\alpha
=
\text{damping parameter},
\]

and

\[
\mathbf{v}
=
\text{teleportation distribution}.
\]

%%%%%%%%%%%%%%%%%%%%%%%%%%%%%%%%%%%%%%%%%%%%%%%%%%%%%%%%%%%%
\subsection{Katz Centrality}
\label{subsec:katz-centrality}
%%%%%%%%%%%%%%%%%%%%%%%%%%%%%%%%%%%%%%%%%%%%%%%%%%%%%%%%%%%%

Katz centrality counts walks of many lengths while reducing the influence of
longer walks.

A common form is

\[
\boxed{
\mathbf{x}
=
\alpha A\mathbf{x}
+
\beta\mathbf{1}.
}
\]

Thus

\[
\boxed{
\mathbf{x}
=
\beta
(I-\alpha A)^{-1}\mathbf{1},
}
\]

when the inverse exists.

Typically,

\[
\alpha
<
\frac{1}{\rho(A)},
\]

where \(\rho(A)\) is the spectral radius.

%%%%%%%%%%%%%%%%%%%%%%%%%%%%%%%%%%%%%%%%%%%%%%%%%%%%%%%%%%%%
\subsection{Graph Laplacian}
\label{subsec:graph-laplacian-network}
%%%%%%%%%%%%%%%%%%%%%%%%%%%%%%%%%%%%%%%%%%%%%%%%%%%%%%%%%%%%

For an undirected graph, let

\[
D
=
\operatorname{diag}(k_1,\ldots,k_n)
\]

be the degree matrix.

The graph Laplacian is

\[
\boxed{
L
=
D-A.
}
\]

It plays a central role in connectivity, diffusion, consensus, and spectral
clustering.

%%%%%%%%%%%%%%%%%%%%%%%%%%%%%%%%%%%%%%%%%%%%%%%%%%%%%%%%%%%%
\subsection{Laplacian Quadratic Form}
\label{subsec:laplacian-quadratic-form}
%%%%%%%%%%%%%%%%%%%%%%%%%%%%%%%%%%%%%%%%%%%%%%%%%%%%%%%%%%%%

For vector \(\mathbf{x}\),

\[
\boxed{
\mathbf{x}^TL\mathbf{x}
=
\frac12
\sum_{i,j}
A_{ij}
(x_i-x_j)^2.
}
\]

For an undirected graph,

\[
\boxed{
\mathbf{x}^TL\mathbf{x}\geq0.
}
\]

Thus \(L\) is positive semidefinite.

%%%%%%%%%%%%%%%%%%%%%%%%%%%%%%%%%%%%%%%%%%%%%%%%%%%%%%%%%%%%
\subsection{Zero Laplacian Eigenvalue}
\label{subsec:zero-laplacian-eigenvalue}
%%%%%%%%%%%%%%%%%%%%%%%%%%%%%%%%%%%%%%%%%%%%%%%%%%%%%%%%%%%%

Since

\[
L\mathbf{1}=0,
\]

we always have eigenvalue

\[
\boxed{
\lambda_1=0.
}
\]

For an undirected graph, the multiplicity of eigenvalue \(0\) equals the
number of connected components.

%%%%%%%%%%%%%%%%%%%%%%%%%%%%%%%%%%%%%%%%%%%%%%%%%%%%%%%%%%%%
\subsection{Algebraic Connectivity}
\label{subsec:algebraic-connectivity}
%%%%%%%%%%%%%%%%%%%%%%%%%%%%%%%%%%%%%%%%%%%%%%%%%%%%%%%%%%%%

For a connected undirected graph, order Laplacian eigenvalues as

\[
0
=
\lambda_1
<
\lambda_2
\leq
\cdots
\leq
\lambda_n.
\]

The value

\[
\boxed{
\lambda_2
}
\]

is called algebraic connectivity.

Larger values generally indicate stronger graph connectivity.

%%%%%%%%%%%%%%%%%%%%%%%%%%%%%%%%%%%%%%%%%%%%%%%%%%%%%%%%%%%%
\subsection{Fiedler Vector}
\label{subsec:fiedler-vector}
%%%%%%%%%%%%%%%%%%%%%%%%%%%%%%%%%%%%%%%%%%%%%%%%%%%%%%%%%%%%

An eigenvector associated with

\[
\lambda_2
\]

is called a Fiedler vector.

Its signs or values can be used to divide a graph into groups.

This provides a foundation for spectral partitioning and clustering.

%%%%%%%%%%%%%%%%%%%%%%%%%%%%%%%%%%%%%%%%%%%%%%%%%%%%%%%%%%%%
\subsection{Normalized Laplacians}
\label{subsec:normalized-laplacian}
%%%%%%%%%%%%%%%%%%%%%%%%%%%%%%%%%%%%%%%%%%%%%%%%%%%%%%%%%%%%

A symmetric normalized Laplacian is

\[
\boxed{
L_{\mathrm{sym}}
=
I
-
D^{-1/2}AD^{-1/2}.
}
\]

Another form is

\[
\boxed{
L_{\mathrm{rw}}
=
I-D^{-1}A.
}
\]

These are useful when node degrees vary strongly.

%%%%%%%%%%%%%%%%%%%%%%%%%%%%%%%%%%%%%%%%%%%%%%%%%%%%%%%%%%%%
\subsection{Spectral Clustering}
\label{subsec:spectral-clustering-network}
%%%%%%%%%%%%%%%%%%%%%%%%%%%%%%%%%%%%%%%%%%%%%%%%%%%%%%%%%%%%

Spectral clustering uses eigenvectors of a Laplacian matrix to embed nodes in
a low-dimensional space.

A clustering algorithm is then applied to the embedded coordinates.

The method converts graph partitioning into a linear-algebra problem.

%%%%%%%%%%%%%%%%%%%%%%%%%%%%%%%%%%%%%%%%%%%%%%%%%%%%%%%%%%%%
\subsection{Random Networks}
\label{subsec:random-networks}
%%%%%%%%%%%%%%%%%%%%%%%%%%%%%%%%%%%%%%%%%%%%%%%%%%%%%%%%%%%%

A random graph is generated by a probabilistic rule.

Random networks provide:

\[
\boxed{
\text{null models},
}
\]

\[
\boxed{
\text{baseline comparisons},
}
\]

and

\[
\boxed{
\text{models of uncertain connectivity}.
}
\]

%%%%%%%%%%%%%%%%%%%%%%%%%%%%%%%%%%%%%%%%%%%%%%%%%%%%%%%%%%%%
\subsection{Erdos--Renyi Model}
\label{subsec:erdos-renyi-network}
%%%%%%%%%%%%%%%%%%%%%%%%%%%%%%%%%%%%%%%%%%%%%%%%%%%%%%%%%%%%

In the \(G(n,p)\) model, there are \(n\) vertices and each possible
undirected edge is independently present with probability \(p\).

There are

\[
\binom n2
\]

possible edges.

The expected number of edges is

\[
\boxed{
\mathbb{E}[|E|]
=
p\binom n2.
}
\]

%%%%%%%%%%%%%%%%%%%%%%%%%%%%%%%%%%%%%%%%%%%%%%%%%%%%%%%%%%%%
\subsection{Expected Degree in a Random Graph}
\label{subsec:random-graph-expected-degree}
%%%%%%%%%%%%%%%%%%%%%%%%%%%%%%%%%%%%%%%%%%%%%%%%%%%%%%%%%%%%

For \(G(n,p)\), each node has \(n-1\) possible neighbors.

Thus

\[
\boxed{
\mathbb{E}[k_i]
=
(n-1)p.
}
\]

The average expected degree is therefore

\[
\boxed{
\langle k\rangle
=
(n-1)p.
}
\]

%%%%%%%%%%%%%%%%%%%%%%%%%%%%%%%%%%%%%%%%%%%%%%%%%%%%%%%%%%%%
\subsection{Random-Graph Degree Distribution}
\label{subsec:random-graph-degree-distribution}
%%%%%%%%%%%%%%%%%%%%%%%%%%%%%%%%%%%%%%%%%%%%%%%%%%%%%%%%%%%%

For \(G(n,p)\),

\[
\boxed{
K_i
\sim
\operatorname{Binomial}(n-1,p).
}
\]

Therefore,

\[
\boxed{
\Pr(K_i=k)
=
\binom{n-1}{k}
p^k
(1-p)^{n-1-k}.
}
\]

In suitable sparse limits, this approaches a Poisson distribution.

%%%%%%%%%%%%%%%%%%%%%%%%%%%%%%%%%%%%%%%%%%%%%%%%%%%%%%%%%%%%
\subsection{Giant Components}
\label{subsec:giant-components}
%%%%%%%%%%%%%%%%%%%%%%%%%%%%%%%%%%%%%%%%%%%%%%%%%%%%%%%%%%%%

Random graphs can undergo a connectivity transition.

When average degree becomes sufficiently large, a connected component
containing a positive fraction of all nodes can emerge.

This is called a giant component.

Such threshold phenomena are central to network percolation.

%%%%%%%%%%%%%%%%%%%%%%%%%%%%%%%%%%%%%%%%%%%%%%%%%%%%%%%%%%%%
\subsection{Small-World Networks}
\label{subsec:small-world-networks}
%%%%%%%%%%%%%%%%%%%%%%%%%%%%%%%%%%%%%%%%%%%%%%%%%%%%%%%%%%%%

Small-world networks combine:

\[
\boxed{
\text{high local clustering}
}
\]

with

\[
\boxed{
\text{short global path lengths}.
}
\]

A small number of long-range links can dramatically reduce network
distances.

%%%%%%%%%%%%%%%%%%%%%%%%%%%%%%%%%%%%%%%%%%%%%%%%%%%%%%%%%%%%
\subsection{Watts--Strogatz Model}
\label{subsec:watts-strogatz-model}
%%%%%%%%%%%%%%%%%%%%%%%%%%%%%%%%%%%%%%%%%%%%%%%%%%%%%%%%%%%%

The Watts--Strogatz construction begins with a regular ring lattice.

Edges are then rewired randomly with some probability.

This interpolates between:

\[
\boxed{
\text{regular structure}
}
\]

and

\[
\boxed{
\text{random structure}.
}
\]

Intermediate rewiring can produce small-world behavior.

%%%%%%%%%%%%%%%%%%%%%%%%%%%%%%%%%%%%%%%%%%%%%%%%%%%%%%%%%%%%
\subsection{Preferential Attachment}
\label{subsec:preferential-attachment}
%%%%%%%%%%%%%%%%%%%%%%%%%%%%%%%%%%%%%%%%%%%%%%%%%%%%%%%%%%%%

Preferential attachment models growing networks in which new nodes are more
likely to connect to already well-connected nodes.

Schematically,

\[
\boxed{
\Pr(\text{new node attaches to }i)
\propto
k_i.
}
\]

This is sometimes summarized as:

\[
\boxed{
\text{the rich get richer}.
}
\]

%%%%%%%%%%%%%%%%%%%%%%%%%%%%%%%%%%%%%%%%%%%%%%%%%%%%%%%%%%%%
\subsection{Scale-Free Networks}
\label{subsec:scale-free-networks}
%%%%%%%%%%%%%%%%%%%%%%%%%%%%%%%%%%%%%%%%%%%%%%%%%%%%%%%%%%%%

Some networks exhibit broad degree distributions approximated over a range by

\[
\boxed{
P(k)
\propto
k^{-\gamma}.
}
\]

Such networks are often described as scale-free.

However, observing a heavy tail does not by itself prove a power-law degree
distribution.

Statistical testing is required.

%%%%%%%%%%%%%%%%%%%%%%%%%%%%%%%%%%%%%%%%%%%%%%%%%%%%%%%%%%%%
\subsection{Configuration Model}
\label{subsec:configuration-model}
%%%%%%%%%%%%%%%%%%%%%%%%%%%%%%%%%%%%%%%%%%%%%%%%%%%%%%%%%%%%

The configuration model generates random networks with a specified degree
sequence.

Each node is assigned a number of stubs equal to its degree.

Stubs are paired to form edges.

Depending on the version, self-loops or repeated edges may occur.

%%%%%%%%%%%%%%%%%%%%%%%%%%%%%%%%%%%%%%%%%%%%%%%%%%%%%%%%%%%%
\subsection{Assortativity}
\label{subsec:network-assortativity}
%%%%%%%%%%%%%%%%%%%%%%%%%%%%%%%%%%%%%%%%%%%%%%%%%%%%%%%%%%%%

A network is assortative with respect to an attribute if similar nodes tend
to connect.

Degree assortativity measures whether high-degree nodes preferentially
connect to other high-degree nodes.

Positive assortativity indicates similarity.

Negative assortativity indicates disassortative mixing.

%%%%%%%%%%%%%%%%%%%%%%%%%%%%%%%%%%%%%%%%%%%%%%%%%%%%%%%%%%%%
\subsection{Homophily}
\label{subsec:network-homophily}
%%%%%%%%%%%%%%%%%%%%%%%%%%%%%%%%%%%%%%%%%%%%%%%%%%%%%%%%%%%%

Homophily describes a tendency for nodes with similar characteristics to
connect.

Examples include similarity in:

\[
\boxed{
\text{location},
}
\]

\[
\boxed{
\text{age},
}
\]

\[
\boxed{
\text{interests},
}
\]

or other attributes.

Observed similarity can arise from both selection and network influence, so
causal interpretation requires care.

%%%%%%%%%%%%%%%%%%%%%%%%%%%%%%%%%%%%%%%%%%%%%%%%%%%%%%%%%%%%
\subsection{Community Structure}
\label{subsec:community-structure}
%%%%%%%%%%%%%%%%%%%%%%%%%%%%%%%%%%%%%%%%%%%%%%%%%%%%%%%%%%%%

A community is a group of nodes with relatively dense internal connections
and weaker connections to the rest of the network.

Community structure can reveal:

\[
\boxed{
\text{functional groups},
}
\]

\[
\boxed{
\text{social groups},
}
\]

or

\[
\boxed{
\text{organizational modules}.
}
\]

There is no single universally correct mathematical definition of a
community.

%%%%%%%%%%%%%%%%%%%%%%%%%%%%%%%%%%%%%%%%%%%%%%%%%%%%%%%%%%%%
\subsection{Modularity}
\label{subsec:network-modularity}
%%%%%%%%%%%%%%%%%%%%%%%%%%%%%%%%%%%%%%%%%%%%%%%%%%%%%%%%%%%%

Modularity compares the observed number of within-community edges with a
null-model expectation.

One common form is

\[
\boxed{
Q
=
\frac{1}{2m}
\sum_{i,j}
\left(
A_{ij}
-
\frac{k_ik_j}{2m}
\right)
\delta(c_i,c_j),
}
\]

where

\[
m=|E|.
\]

Larger \(Q\) indicates more within-group connectivity than expected under the
chosen null model.

%%%%%%%%%%%%%%%%%%%%%%%%%%%%%%%%%%%%%%%%%%%%%%%%%%%%%%%%%%%%
\subsection{Community Detection}
\label{subsec:community-detection}
%%%%%%%%%%%%%%%%%%%%%%%%%%%%%%%%%%%%%%%%%%%%%%%%%%%%%%%%%%%%

Community detection methods include:

\[
\boxed{
\text{modularity optimization},
}
\]

\[
\boxed{
\text{spectral methods},
}
\]

\[
\boxed{
\text{hierarchical clustering},
}
\]

and

\[
\boxed{
\text{probabilistic block models}.
}
\]

Different methods may return different partitions.

%%%%%%%%%%%%%%%%%%%%%%%%%%%%%%%%%%%%%%%%%%%%%%%%%%%%%%%%%%%%
\subsection{Stochastic Block Models}
\label{subsec:stochastic-block-models}
%%%%%%%%%%%%%%%%%%%%%%%%%%%%%%%%%%%%%%%%%%%%%%%%%%%%%%%%%%%%

A stochastic block model assigns nodes to latent groups.

The probability of an edge depends on the groups of its endpoints.

If node \(i\) belongs to group \(g_i\), then

\[
\boxed{
\Pr(A_{ij}=1)
=
P_{g_i g_j}.
}
\]

This provides a statistical model of community structure.

%%%%%%%%%%%%%%%%%%%%%%%%%%%%%%%%%%%%%%%%%%%%%%%%%%%%%%%%%%%%
\subsection{Hierarchical Networks}
\label{subsec:hierarchical-networks}
%%%%%%%%%%%%%%%%%%%%%%%%%%%%%%%%%%%%%%%%%%%%%%%%%%%%%%%%%%%%

Some networks contain nested structure.

Small communities may combine into larger communities, which themselves form
larger-scale groups.

Hierarchical clustering represents this structure using a tree or
dendrogram.

%%%%%%%%%%%%%%%%%%%%%%%%%%%%%%%%%%%%%%%%%%%%%%%%%%%%%%%%%%%%
\subsection{Bipartite Networks}
\label{subsec:bipartite-networks}
%%%%%%%%%%%%%%%%%%%%%%%%%%%%%%%%%%%%%%%%%%%%%%%%%%%%%%%%%%%%

A bipartite network has node sets

\[
\boxed{
U
}
\]

and

\[
\boxed{
V
}
\]

with edges only between the sets.

Examples include:

\[
\boxed{
\text{users and products},
}
\]

\[
\boxed{
\text{authors and papers},
}
\]

and

\[
\boxed{
\text{pollinators and plants}.
}
\]

%%%%%%%%%%%%%%%%%%%%%%%%%%%%%%%%%%%%%%%%%%%%%%%%%%%%%%%%%%%%
\subsection{Bipartite Projection}
\label{subsec:bipartite-projection}
%%%%%%%%%%%%%%%%%%%%%%%%%%%%%%%%%%%%%%%%%%%%%%%%%%%%%%%%%%%%

A bipartite network can be projected onto one node set.

Two nodes in \(U\) are connected if they share a neighbor in \(V\).

Projection can create dense networks and may lose information, so the
original bipartite structure should often be retained.

%%%%%%%%%%%%%%%%%%%%%%%%%%%%%%%%%%%%%%%%%%%%%%%%%%%%%%%%%%%%
\subsection{Multilayer Networks}
\label{subsec:multilayer-networks}
%%%%%%%%%%%%%%%%%%%%%%%%%%%%%%%%%%%%%%%%%%%%%%%%%%%%%%%%%%%%

A multilayer network contains multiple types of relationships.

For example, people may interact through:

\[
\boxed{
\text{work},
}
\]

\[
\boxed{
\text{family},
}
\]

and

\[
\boxed{
\text{online communication}.
}
\]

Each relation can be represented as a different layer.

%%%%%%%%%%%%%%%%%%%%%%%%%%%%%%%%%%%%%%%%%%%%%%%%%%%%%%%%%%%%
\subsection{Multiplex Networks}
\label{subsec:multiplex-networks}
%%%%%%%%%%%%%%%%%%%%%%%%%%%%%%%%%%%%%%%%%%%%%%%%%%%%%%%%%%%%

A multiplex network is a multilayer system in which the same nodes appear in
multiple layers.

Different edge types connect the same collection of entities.

Dynamics can depend on interactions across layers.

%%%%%%%%%%%%%%%%%%%%%%%%%%%%%%%%%%%%%%%%%%%%%%%%%%%%%%%%%%%%
\subsection{Temporal Networks}
\label{subsec:temporal-networks}
%%%%%%%%%%%%%%%%%%%%%%%%%%%%%%%%%%%%%%%%%%%%%%%%%%%%%%%%%%%%

In a temporal network,

\[
\boxed{
E=E(t).
}
\]

Connections can appear and disappear over time.

Temporal ordering matters because a sequence of edges may permit a
time-respecting path even when the static aggregate network is misleading.

%%%%%%%%%%%%%%%%%%%%%%%%%%%%%%%%%%%%%%%%%%%%%%%%%%%%%%%%%%%%
\subsection{Network Robustness}
\label{subsec:network-robustness}
%%%%%%%%%%%%%%%%%%%%%%%%%%%%%%%%%%%%%%%%%%%%%%%%%%%%%%%%%%%%

Network robustness studies how connectivity and function change under node
or edge removal.

Failures may be:

\[
\boxed{
\text{random}
}
\]

or

\[
\boxed{
\text{targeted}.
}
\]

Different network structures respond differently to these disruptions.

%%%%%%%%%%%%%%%%%%%%%%%%%%%%%%%%%%%%%%%%%%%%%%%%%%%%%%%%%%%%
\subsection{Site Percolation}
\label{subsec:site-percolation-network}
%%%%%%%%%%%%%%%%%%%%%%%%%%%%%%%%%%%%%%%%%%%%%%%%%%%%%%%%%%%%

In site percolation, nodes are independently retained or removed according to
some probability.

One studies whether a large connected component survives.

This models random node failures.

%%%%%%%%%%%%%%%%%%%%%%%%%%%%%%%%%%%%%%%%%%%%%%%%%%%%%%%%%%%%
\subsection{Bond Percolation}
\label{subsec:bond-percolation-network}
%%%%%%%%%%%%%%%%%%%%%%%%%%%%%%%%%%%%%%%%%%%%%%%%%%%%%%%%%%%%

In bond percolation, edges are independently retained or removed.

This models failure or activation of connections.

Percolation theory studies threshold behavior in global connectivity.

%%%%%%%%%%%%%%%%%%%%%%%%%%%%%%%%%%%%%%%%%%%%%%%%%%%%%%%%%%%%
\subsection{Cascading Failures}
\label{subsec:cascading-failures}
%%%%%%%%%%%%%%%%%%%%%%%%%%%%%%%%%%%%%%%%%%%%%%%%%%%%%%%%%%%%

A failure may increase load on neighboring or alternative components.

Those components may then fail, producing a cascade.

Examples include:

\[
\boxed{
\text{power grids},
}
\]

\[
\boxed{
\text{financial systems},
}
\]

and

\[
\boxed{
\text{transportation networks}.
}
\]

%%%%%%%%%%%%%%%%%%%%%%%%%%%%%%%%%%%%%%%%%%%%%%%%%%%%%%%%%%%%
\subsection{Random Walks on Networks}
\label{subsec:random-walks-networks}
%%%%%%%%%%%%%%%%%%%%%%%%%%%%%%%%%%%%%%%%%%%%%%%%%%%%%%%%%%%%

For a simple random walk on an undirected graph, a walker at node \(i\)
chooses one of its \(k_i\) neighbors uniformly.

The transition probability is

\[
\boxed{
P_{ij}
=
\frac{A_{ij}}
{k_i},
}
\]

when \(k_i>0\).

Thus

\[
\boxed{
P=D^{-1}A.
}
\]

%%%%%%%%%%%%%%%%%%%%%%%%%%%%%%%%%%%%%%%%%%%%%%%%%%%%%%%%%%%%
\subsection{Probability Evolution}
\label{subsec:network-random-walk-evolution}
%%%%%%%%%%%%%%%%%%%%%%%%%%%%%%%%%%%%%%%%%%%%%%%%%%%%%%%%%%%%

Let row vector

\[
\mathbf{p}_t
\]

give node probabilities at step \(t\).

Then

\[
\boxed{
\mathbf{p}_{t+1}
=
\mathbf{p}_tP.
}
\]

This is a Markov chain on the network.

%%%%%%%%%%%%%%%%%%%%%%%%%%%%%%%%%%%%%%%%%%%%%%%%%%%%%%%%%%%%
\subsection{Stationary Distribution}
\label{subsec:network-stationary-distribution}
%%%%%%%%%%%%%%%%%%%%%%%%%%%%%%%%%%%%%%%%%%%%%%%%%%%%%%%%%%%%

For a connected non-bipartite undirected graph, the simple random walk has
stationary distribution

\[
\boxed{
\pi_i
=
\frac{k_i}
{\sum_jk_j}.
}
\]

Using the handshaking relation,

\[
\boxed{
\pi_i
=
\frac{k_i}
{2|E|}.
}
\]

Higher-degree nodes receive greater stationary probability.

%%%%%%%%%%%%%%%%%%%%%%%%%%%%%%%%%%%%%%%%%%%%%%%%%%%%%%%%%%%%
\subsection{Worked Example: Random-Walk Stationary Distribution}
\label{subsec:worked-random-walk-stationary}
%%%%%%%%%%%%%%%%%%%%%%%%%%%%%%%%%%%%%%%%%%%%%%%%%%%%%%%%%%%%

\begin{workedexamplebox}{Path of Three Nodes}

For the path

\[
1-2-3,
\]

degrees are

\[
k_1=1,
\qquad
k_2=2,
\qquad
k_3=1.
\]

Their sum is

\[
4.
\]

Therefore,

\[
\pi_1=\frac14,
\qquad
\pi_2=\frac12,
\qquad
\pi_3=\frac14.
\]

\begin{answerbox}
\[
\boxed{
\boldsymbol{\pi}
=
\left(
\frac14,
\frac12,
\frac14
\right).
}
\]
\end{answerbox}

\end{workedexamplebox}

%%%%%%%%%%%%%%%%%%%%%%%%%%%%%%%%%%%%%%%%%%%%%%%%%%%%%%%%%%%%
\subsection{Diffusion on Networks}
\label{subsec:diffusion-on-networks}
%%%%%%%%%%%%%%%%%%%%%%%%%%%%%%%%%%%%%%%%%%%%%%%%%%%%%%%%%%%%

Let \(x_i(t)\) be a scalar value on node \(i\).

Diffusion may be modeled as

\[
\boxed{
\dot{x}_i
=
\kappa
\sum_j
A_{ij}
(x_j-x_i).
}
\]

Using the graph Laplacian,

\[
\boxed{
\dot{\mathbf{x}}
=
-\kappa L\mathbf{x}.
}
\]

%%%%%%%%%%%%%%%%%%%%%%%%%%%%%%%%%%%%%%%%%%%%%%%%%%%%%%%%%%%%
\subsection{Solution of Network Diffusion}
\label{subsec:network-diffusion-solution}
%%%%%%%%%%%%%%%%%%%%%%%%%%%%%%%%%%%%%%%%%%%%%%%%%%%%%%%%%%%%

The solution is

\[
\boxed{
\mathbf{x}(t)
=
e^{-\kappa Lt}
\mathbf{x}(0).
}
\]

If the graph is connected, differences among node values decay under this
simple diffusion process.

The zero Laplacian mode represents the conserved uniform component.

%%%%%%%%%%%%%%%%%%%%%%%%%%%%%%%%%%%%%%%%%%%%%%%%%%%%%%%%%%%%
\subsection{Consensus}
\label{subsec:network-consensus}
%%%%%%%%%%%%%%%%%%%%%%%%%%%%%%%%%%%%%%%%%%%%%%%%%%%%%%%%%%%%

Consensus dynamics are

\[
\boxed{
\dot{\mathbf{x}}
=
-L\mathbf{x}.
}
\]

For a connected undirected graph,

\[
\mathbf{x}(t)
\]

approaches a common value.

For this symmetric consensus model, the limiting value is the initial
average:

\[
\boxed{
x_i(t)
\to
\frac1n
\sum_{j=1}^{n}x_j(0).
}
\]

%%%%%%%%%%%%%%%%%%%%%%%%%%%%%%%%%%%%%%%%%%%%%%%%%%%%%%%%%%%%
\subsection{Algebraic Connectivity and Consensus Speed}
\label{subsec:algebraic-connectivity-consensus}
%%%%%%%%%%%%%%%%%%%%%%%%%%%%%%%%%%%%%%%%%%%%%%%%%%%%%%%%%%%%

The slowest nonconstant decay mode is governed by

\[
\lambda_2.
\]

Thus larger algebraic connectivity typically produces faster convergence of
continuous-time consensus:

\[
\boxed{
\text{larger }\lambda_2
\Rightarrow
\text{faster mixing}.
}
\]

%%%%%%%%%%%%%%%%%%%%%%%%%%%%%%%%%%%%%%%%%%%%%%%%%%%%%%%%%%%%
\subsection{Synchronization}
\label{subsec:network-synchronization}
%%%%%%%%%%%%%%%%%%%%%%%%%%%%%%%%%%%%%%%%%%%%%%%%%%%%%%%%%%%%

Synchronization occurs when coupled dynamical systems evolve coherently.

A generic model is

\[
\boxed{
\dot{\mathbf{x}}_i
=
f(\mathbf{x}_i)
+
\kappa
\sum_j
A_{ij}
H(\mathbf{x}_j-\mathbf{x}_i).
}
\]

Network topology can determine whether synchronized behavior is stable.

%%%%%%%%%%%%%%%%%%%%%%%%%%%%%%%%%%%%%%%%%%%%%%%%%%%%%%%%%%%%
\subsection{Kuramoto Model}
\label{subsec:kuramoto-model}
%%%%%%%%%%%%%%%%%%%%%%%%%%%%%%%%%%%%%%%%%%%%%%%%%%%%%%%%%%%%

A classical phase-oscillator model is

\[
\boxed{
\dot{\theta}_i
=
\omega_i
+
\frac{K}{n}
\sum_{j=1}^{n}
A_{ij}
\sin(\theta_j-\theta_i).
}
\]

Here:

\[
\theta_i
=
\text{phase},
\]

\[
\omega_i
=
\text{natural frequency},
\]

and

\[
K
=
\text{coupling strength}.
\]

Increasing coupling can promote collective synchronization.

%%%%%%%%%%%%%%%%%%%%%%%%%%%%%%%%%%%%%%%%%%%%%%%%%%%%%%%%%%%%
\subsection{Network Epidemics}
\label{subsec:network-epidemics}
%%%%%%%%%%%%%%%%%%%%%%%%%%%%%%%%%%%%%%%%%%%%%%%%%%%%%%%%%%%%

In network epidemic models, infection is transmitted along edges rather than
through homogeneous mixing.

The infection risk of node \(i\) depends on the infection states of its
neighbors.

A conceptual infection pressure is

\[
\boxed{
\beta
\sum_j
A_{ij}I_j.
}
\]

Network topology can strongly affect epidemic thresholds and outbreak size.

%%%%%%%%%%%%%%%%%%%%%%%%%%%%%%%%%%%%%%%%%%%%%%%%%%%%%%%%%%%%
\subsection{Spectral Epidemic Thresholds}
\label{subsec:spectral-epidemic-thresholds}
%%%%%%%%%%%%%%%%%%%%%%%%%%%%%%%%%%%%%%%%%%%%%%%%%%%%%%%%%%%%

In several network epidemic approximations, epidemic persistence is related
to the spectral radius

\[
\boxed{
\rho(A).
}
\]

A larger dominant adjacency eigenvalue generally increases spreading
potential.

The precise threshold depends on the epidemic model and approximation being
used.

%%%%%%%%%%%%%%%%%%%%%%%%%%%%%%%%%%%%%%%%%%%%%%%%%%%%%%%%%%%%
\subsection{Influence and Contagion}
\label{subsec:network-contagion}
%%%%%%%%%%%%%%%%%%%%%%%%%%%%%%%%%%%%%%%%%%%%%%%%%%%%%%%%%%%%

Not all spreading processes behave like infectious disease.

Social contagion may require:

\[
\boxed{
\text{one exposure}
}
\]

or

\[
\boxed{
\text{multiple reinforcing exposures}.
}
\]

Simple contagion and complex contagion can therefore respond differently to
clustering and community structure.

%%%%%%%%%%%%%%%%%%%%%%%%%%%%%%%%%%%%%%%%%%%%%%%%%%%%%%%%%%%%
\subsection{Threshold Models}
\label{subsec:network-threshold-models}
%%%%%%%%%%%%%%%%%%%%%%%%%%%%%%%%%%%%%%%%%%%%%%%%%%%%%%%%%%%%

Suppose node \(i\) adopts a state when the fraction of active neighbors
exceeds threshold \(\theta_i\).

Then adoption occurs when

\[
\boxed{
\frac{
\sum_jA_{ij}x_j
}{
k_i
}
\geq
\theta_i.
}
\]

Local thresholds can generate global cascades.

%%%%%%%%%%%%%%%%%%%%%%%%%%%%%%%%%%%%%%%%%%%%%%%%%%%%%%%%%%%%
\subsection{Flow Networks}
\label{subsec:network-flow-networks}
%%%%%%%%%%%%%%%%%%%%%%%%%%%%%%%%%%%%%%%%%%%%%%%%%%%%%%%%%%%%

A flow network is a directed graph with edge capacities

\[
\boxed{
c_{ij}\geq0.
}
\]

A feasible flow satisfies capacity constraints

\[
\boxed{
0\leq f_{ij}\leq c_{ij}
}
\]

and conservation at intermediate nodes.

%%%%%%%%%%%%%%%%%%%%%%%%%%%%%%%%%%%%%%%%%%%%%%%%%%%%%%%%%%%%
\subsection{Flow Conservation}
\label{subsec:network-flow-conservation}
%%%%%%%%%%%%%%%%%%%%%%%%%%%%%%%%%%%%%%%%%%%%%%%%%%%%%%%%%%%%

For a node \(i\) other than source or sink,

\[
\boxed{
\sum_jf_{ji}
=
\sum_jf_{ij}.
}
\]

Thus total inflow equals total outflow.

This is a discrete conservation law.

%%%%%%%%%%%%%%%%%%%%%%%%%%%%%%%%%%%%%%%%%%%%%%%%%%%%%%%%%%%%
\subsection{Maximum Flow}
\label{subsec:network-maximum-flow}
%%%%%%%%%%%%%%%%%%%%%%%%%%%%%%%%%%%%%%%%%%%%%%%%%%%%%%%%%%%%

The maximum-flow problem seeks the largest feasible flow from a source \(s\)
to a sink \(t\).

Its optimization structure is

\[
\boxed{
\max
\quad
\text{flow value}
}
\]

subject to capacity and conservation constraints.

%%%%%%%%%%%%%%%%%%%%%%%%%%%%%%%%%%%%%%%%%%%%%%%%%%%%%%%%%%%%
\subsection{Minimum Cut}
\label{subsec:network-minimum-cut}
%%%%%%%%%%%%%%%%%%%%%%%%%%%%%%%%%%%%%%%%%%%%%%%%%%%%%%%%%%%%

An \(s\)-\(t\) cut partitions the vertices into sets \(S\) and \(T\) with

\[
s\in S,
\qquad
t\in T.
\]

Its capacity is the total capacity of directed edges from \(S\) to \(T\).

The max-flow min-cut theorem states

\[
\boxed{
\text{maximum flow value}
=
\text{minimum cut capacity}.
}
\]

%%%%%%%%%%%%%%%%%%%%%%%%%%%%%%%%%%%%%%%%%%%%%%%%%%%%%%%%%%%%
\subsection{Shortest-Path Networks}
\label{subsec:network-shortest-paths}
%%%%%%%%%%%%%%%%%%%%%%%%%%%%%%%%%%%%%%%%%%%%%%%%%%%%%%%%%%%%

When edge weights represent costs or distances, route optimization seeks

\[
\boxed{
\min
\sum_{(i,j)\in P}
w_{ij}
}
\]

over paths \(P\) connecting specified nodes.

Shortest paths are central in:

\[
\boxed{
\text{transportation},
}
\]

\[
\boxed{
\text{communication},
}
\]

and

\[
\boxed{
\text{logistics}.
}
\]

%%%%%%%%%%%%%%%%%%%%%%%%%%%%%%%%%%%%%%%%%%%%%%%%%%%%%%%%%%%%
\subsection{Network Resilience}
\label{subsec:network-resilience}
%%%%%%%%%%%%%%%%%%%%%%%%%%%%%%%%%%%%%%%%%%%%%%%%%%%%%%%%%%%%

Resilience refers to a network's ability to retain or recover function after
disruption.

Possible metrics include:

\[
\boxed{
\text{component size},
}
\]

\[
\boxed{
\text{path efficiency},
}
\]

\[
\boxed{
\text{flow capacity},
}
\]

and

\[
\boxed{
\text{recovery time}.
}
\]

Structural connectivity alone does not completely determine functional
resilience.

%%%%%%%%%%%%%%%%%%%%%%%%%%%%%%%%%%%%%%%%%%%%%%%%%%%%%%%%%%%%
\subsection{Social Networks}
\label{subsec:social-networks}
%%%%%%%%%%%%%%%%%%%%%%%%%%%%%%%%%%%%%%%%%%%%%%%%%%%%%%%%%%%%

Nodes may represent individuals or organizations.

Edges may represent:

\[
\boxed{
\text{friendship},
}
\]

\[
\boxed{
\text{communication},
}
\]

\[
\boxed{
\text{collaboration},
}
\]

or

\[
\boxed{
\text{information flow}.
}
\]

Centrality, communities, homophily, and diffusion are common topics.

%%%%%%%%%%%%%%%%%%%%%%%%%%%%%%%%%%%%%%%%%%%%%%%%%%%%%%%%%%%%
\subsection{Biological Networks}
\label{subsec:network-biological-networks}
%%%%%%%%%%%%%%%%%%%%%%%%%%%%%%%%%%%%%%%%%%%%%%%%%%%%%%%%%%%%

Biological networks include:

\[
\boxed{
\text{gene-regulatory networks},
}
\]

\[
\boxed{
\text{protein-interaction networks},
}
\]

\[
\boxed{
\text{metabolic networks},
}
\]

and

\[
\boxed{
\text{neural networks}.
}
\]

Network topology can influence biological robustness and dynamics.

%%%%%%%%%%%%%%%%%%%%%%%%%%%%%%%%%%%%%%%%%%%%%%%%%%%%%%%%%%%%
\subsection{Ecological Networks}
\label{subsec:ecological-networks}
%%%%%%%%%%%%%%%%%%%%%%%%%%%%%%%%%%%%%%%%%%%%%%%%%%%%%%%%%%%%

Food webs connect species through feeding interactions.

Mutualistic networks can connect species such as plants and pollinators.

These networks may be:

\[
\boxed{
\text{directed},
}
\]

\[
\boxed{
\text{weighted},
}
\]

or

\[
\boxed{
\text{bipartite}.
}
\]

%%%%%%%%%%%%%%%%%%%%%%%%%%%%%%%%%%%%%%%%%%%%%%%%%%%%%%%%%%%%
\subsection{Financial Networks}
\label{subsec:financial-networks}
%%%%%%%%%%%%%%%%%%%%%%%%%%%%%%%%%%%%%%%%%%%%%%%%%%%%%%%%%%%%

Financial institutions can be connected through:

\[
\boxed{
\text{loans},
}
\]

\[
\boxed{
\text{shared assets},
}
\]

\[
\boxed{
\text{payment obligations},
}
\]

and

\[
\boxed{
\text{counterparty exposure}.
}
\]

Network structure can affect systemic risk and contagion.

%%%%%%%%%%%%%%%%%%%%%%%%%%%%%%%%%%%%%%%%%%%%%%%%%%%%%%%%%%%%
\subsection{Infrastructure Networks}
\label{subsec:infrastructure-networks}
%%%%%%%%%%%%%%%%%%%%%%%%%%%%%%%%%%%%%%%%%%%%%%%%%%%%%%%%%%%%

Examples include:

\[
\boxed{
\text{power grids},
}
\]

\[
\boxed{
\text{road networks},
}
\]

\[
\boxed{
\text{air transportation},
}
\]

\[
\boxed{
\text{water systems},
}
\]

and

\[
\boxed{
\text{communication networks}.
}
\]

Capacity, reliability, routing, and cascading failure are key modeling issues.

%%%%%%%%%%%%%%%%%%%%%%%%%%%%%%%%%%%%%%%%%%%%%%%%%%%%%%%%%%%%
\subsection{Interdependent Networks}
\label{subsec:interdependent-networks}
%%%%%%%%%%%%%%%%%%%%%%%%%%%%%%%%%%%%%%%%%%%%%%%%%%%%%%%%%%%%

Infrastructure networks often depend on one another.

For example,

\[
\boxed{
\text{communications}
\leftrightarrow
\text{electric power}.
}
\]

Failure in one network can therefore trigger failure in another.

Such dependencies can amplify cascading risk.

%%%%%%%%%%%%%%%%%%%%%%%%%%%%%%%%%%%%%%%%%%%%%%%%%%%%%%%%%%%%
\subsection{Network Motifs}
\label{subsec:network-motifs}
%%%%%%%%%%%%%%%%%%%%%%%%%%%%%%%%%%%%%%%%%%%%%%%%%%%%%%%%%%%%

A network motif is a small subgraph pattern that appears more frequently than
expected under a specified null model.

Examples include:

\[
\boxed{
\text{triangles},
}
\]

\[
\boxed{
\text{feed-forward loops},
}
\]

and other small directed patterns.

Motifs can reveal recurring local organization.

%%%%%%%%%%%%%%%%%%%%%%%%%%%%%%%%%%%%%%%%%%%%%%%%%%%%%%%%%%%%
\subsection{Null Models}
\label{subsec:network-null-models}
%%%%%%%%%%%%%%%%%%%%%%%%%%%%%%%%%%%%%%%%%%%%%%%%%%%%%%%%%%%%

A network statistic becomes more meaningful when compared with a baseline.

A null model specifies what structure would be expected in the absence of the
feature being studied.

Possible null models preserve:

\[
\boxed{
\text{network size},
}
\]

\[
\boxed{
\text{edge count},
}
\]

or

\[
\boxed{
\text{degree sequence}.
}
\]

The choice of null model affects the interpretation.

%%%%%%%%%%%%%%%%%%%%%%%%%%%%%%%%%%%%%%%%%%%%%%%%%%%%%%%%%%%%
\subsection{Network Inference}
\label{subsec:network-inference}
%%%%%%%%%%%%%%%%%%%%%%%%%%%%%%%%%%%%%%%%%%%%%%%%%%%%%%%%%%%%

In many applications, the network is not directly observed.

Instead, edges are inferred from data.

Examples include:

\[
\boxed{
\text{correlations},
}
\]

\[
\boxed{
\text{transactions},
}
\]

\[
\boxed{
\text{communications},
}
\]

or

\[
\boxed{
\text{biological measurements}.
}
\]

The inferred network therefore contains statistical uncertainty.

%%%%%%%%%%%%%%%%%%%%%%%%%%%%%%%%%%%%%%%%%%%%%%%%%%%%%%%%%%%%
\subsection{Link Prediction}
\label{subsec:link-prediction}
%%%%%%%%%%%%%%%%%%%%%%%%%%%%%%%%%%%%%%%%%%%%%%%%%%%%%%%%%%%%

Link prediction estimates which missing or future edges are likely.

Possible predictors include:

\[
\boxed{
\text{common neighbors},
}
\]

\[
\boxed{
\text{node similarity},
}
\]

\[
\boxed{
\text{paths},
}
\]

\[
\boxed{
\text{latent embeddings},
}
\]

and

\[
\boxed{
\text{probabilistic network models}.
}
\]

%%%%%%%%%%%%%%%%%%%%%%%%%%%%%%%%%%%%%%%%%%%%%%%%%%%%%%%%%%%%
\subsection{Common-Neighbor Score}
\label{subsec:common-neighbor-score}
%%%%%%%%%%%%%%%%%%%%%%%%%%%%%%%%%%%%%%%%%%%%%%%%%%%%%%%%%%%%

Let \(N(i)\) be the neighbor set of node \(i\).

A basic link score is

\[
\boxed{
s_{ij}
=
|N(i)\cap N(j)|.
}
\]

Nodes sharing many neighbors may be more likely to connect in networks with
strong triadic closure.

%%%%%%%%%%%%%%%%%%%%%%%%%%%%%%%%%%%%%%%%%%%%%%%%%%%%%%%%%%%%
\subsection{Network Sampling}
\label{subsec:network-sampling}
%%%%%%%%%%%%%%%%%%%%%%%%%%%%%%%%%%%%%%%%%%%%%%%%%%%%%%%%%%%%

Large networks may only be partially observed.

Sampling methods include:

\[
\boxed{
\text{uniform node sampling},
}
\]

\[
\boxed{
\text{edge sampling},
}
\]

and

\[
\boxed{
\text{random-walk sampling}.
}
\]

Sampling can distort degree distributions, clustering, and community
structure.

%%%%%%%%%%%%%%%%%%%%%%%%%%%%%%%%%%%%%%%%%%%%%%%%%%%%%%%%%%%%
\subsection{Missing Edges}
\label{subsec:network-missing-edges}
%%%%%%%%%%%%%%%%%%%%%%%%%%%%%%%%%%%%%%%%%%%%%%%%%%%%%%%%%%%%

An absent observed edge can mean:

\[
\boxed{
\text{no relationship exists}
}
\]

or

\[
\boxed{
\text{the relationship was not observed}.
}
\]

These are different possibilities.

Network-data uncertainty should therefore be incorporated into conclusions.

%%%%%%%%%%%%%%%%%%%%%%%%%%%%%%%%%%%%%%%%%%%%%%%%%%%%%%%%%%%%
\subsection{Weighted Versus Unweighted Analysis}
\label{subsec:weighted-unweighted-analysis}
%%%%%%%%%%%%%%%%%%%%%%%%%%%%%%%%%%%%%%%%%%%%%%%%%%%%%%%%%%%%

Converting a weighted network to an unweighted network discards magnitude
information.

Conversely, using every small measured weight may include noise.

Threshold selection can strongly alter network statistics.

The choice should be justified by the application.

%%%%%%%%%%%%%%%%%%%%%%%%%%%%%%%%%%%%%%%%%%%%%%%%%%%%%%%%%%%%
\subsection{Correlation Networks}
\label{subsec:correlation-networks}
%%%%%%%%%%%%%%%%%%%%%%%%%%%%%%%%%%%%%%%%%%%%%%%%%%%%%%%%%%%%

A correlation network may connect variables when

\[
\boxed{
|\rho_{ij}|
>
\tau,
}
\]

where \(\tau\) is a chosen threshold.

However,

\[
\boxed{
\text{correlation}
\neq
\text{direct interaction}.
}
\]

Common causes and indirect relationships can create correlations.

%%%%%%%%%%%%%%%%%%%%%%%%%%%%%%%%%%%%%%%%%%%%%%%%%%%%%%%%%%%%
\subsection{Network Entropy}
\label{subsec:network-entropy}
%%%%%%%%%%%%%%%%%%%%%%%%%%%%%%%%%%%%%%%%%%%%%%%%%%%%%%%%%%%%

Entropy concepts can quantify uncertainty or heterogeneity in network
structure.

For a normalized degree distribution \(P(k)\), one may define

\[
\boxed{
H
=
-
\sum_k
P(k)\log P(k).
}
\]

This measures uncertainty in the degree of a randomly selected node.

Other forms of network entropy are also used.

%%%%%%%%%%%%%%%%%%%%%%%%%%%%%%%%%%%%%%%%%%%%%%%%%%%%%%%%%%%%
\subsection{Dynamic Networks}
\label{subsec:dynamic-networks}
%%%%%%%%%%%%%%%%%%%%%%%%%%%%%%%%%%%%%%%%%%%%%%%%%%%%%%%%%%%%

Network dynamics can mean two distinct things:

\[
\boxed{
\text{dynamics on a network}
}
\]

or

\[
\boxed{
\text{dynamics of the network itself}.
}
\]

In the first, node states evolve while edges are fixed.

In the second, the edge set changes with time.

Many systems contain both processes simultaneously.

%%%%%%%%%%%%%%%%%%%%%%%%%%%%%%%%%%%%%%%%%%%%%%%%%%%%%%%%%%%%
\subsection{Adaptive Networks}
\label{subsec:adaptive-networks}
%%%%%%%%%%%%%%%%%%%%%%%%%%%%%%%%%%%%%%%%%%%%%%%%%%%%%%%%%%%%

An adaptive network couples state dynamics with structural evolution.

Schematically,

\[
\boxed{
\dot{\mathbf{x}}
=
F(\mathbf{x},A),
}
\]

\[
\boxed{
\dot{A}
=
G(A,\mathbf{x}).
}
\]

Node states affect links, while links affect node states.

This feedback can generate complex collective behavior.

%%%%%%%%%%%%%%%%%%%%%%%%%%%%%%%%%%%%%%%%%%%%%%%%%%%%%%%%%%%%
\subsection{Network Controllability}
\label{subsec:network-controllability}
%%%%%%%%%%%%%%%%%%%%%%%%%%%%%%%%%%%%%%%%%%%%%%%%%%%%%%%%%%%%

Suppose network dynamics are linear:

\[
\boxed{
\dot{\mathbf{x}}
=
A\mathbf{x}
+
B\mathbf{u}.
}
\]

Controllability depends on both:

\[
\boxed{
\text{network dynamics }A
}
\]

and

\[
\boxed{
\text{actuator placement }B.
}
\]

Graph structure alone does not completely determine practical control
difficulty.

%%%%%%%%%%%%%%%%%%%%%%%%%%%%%%%%%%%%%%%%%%%%%%%%%%%%%%%%%%%%
\subsection{Network Observability}
\label{subsec:network-observability}
%%%%%%%%%%%%%%%%%%%%%%%%%%%%%%%%%%%%%%%%%%%%%%%%%%%%%%%%%%%%

For

\[
\dot{\mathbf{x}}
=
A\mathbf{x},
\qquad
\mathbf{y}
=
C\mathbf{x},
\]

sensor placement determines which network states can be inferred.

This connects network science with the observability concepts developed in
Section~\ref{sec:control-theory}.

%%%%%%%%%%%%%%%%%%%%%%%%%%%%%%%%%%%%%%%%%%%%%%%%%%%%%%%%%%%%
\subsection{Optimization on Networks}
\label{subsec:optimization-on-networks}
%%%%%%%%%%%%%%%%%%%%%%%%%%%%%%%%%%%%%%%%%%%%%%%%%%%%%%%%%%%%

Network optimization includes:

\[
\boxed{
\text{shortest paths},
}
\]

\[
\boxed{
\text{minimum spanning trees},
}
\]

\[
\boxed{
\text{maximum flows},
}
\]

\[
\boxed{
\text{routing},
}
\]

\[
\boxed{
\text{facility placement},
}
\]

and

\[
\boxed{
\text{network design}.
}
\]

These topics connect network science with combinatorial optimization and
operations research.

%%%%%%%%%%%%%%%%%%%%%%%%%%%%%%%%%%%%%%%%%%%%%%%%%%%%%%%%%%%%
\subsection{Network Science Workflow}
\label{subsec:network-science-workflow}
%%%%%%%%%%%%%%%%%%%%%%%%%%%%%%%%%%%%%%%%%%%%%%%%%%%%%%%%%%%%

A useful network-analysis workflow is:

\begin{enumerate}
    \item identify the entities represented by nodes;
    \item define precisely what an edge means;
    \item determine whether edges are directed;
    \item determine whether edges are weighted;
    \item determine whether the network changes over time;
    \item identify missing-data and sampling limitations;
    \item construct the adjacency or weighted matrix;
    \item compute basic size and density statistics;
    \item examine degree or strength distributions;
    \item identify connected components;
    \item compute path-based statistics;
    \item examine clustering and triangles;
    \item choose centrality measures appropriate to the question;
    \item examine spectral properties when useful;
    \item test for communities or modular organization;
    \item compare statistics with appropriate null models;
    \item study robustness to node and edge removal;
    \item model relevant dynamical processes;
    \item include direction and weights when physically meaningful;
    \item incorporate temporal or multilayer structure when needed;
    \item quantify uncertainty in inferred networks;
    \item validate interpretations with external information;
    \item avoid treating graph structure alone as proof of causality;
    \item communicate the meaning of nodes, edges, and assumptions clearly.
\end{enumerate}

%%%%%%%%%%%%%%%%%%%%%%%%%%%%%%%%%%%%%%%%%%%%%%%%%%%%%%%%%%%%
\subsection{Common Errors}
\label{subsec:network-science-common-errors}
%%%%%%%%%%%%%%%%%%%%%%%%%%%%%%%%%%%%%%%%%%%%%%%%%%%%%%%%%%%%

\subsubsection{Treating Every Relationship as Undirected}

Many relationships have direction.

Examples include:

\[
\boxed{
\text{citations},
}
\]

\[
\boxed{
\text{payments},
}
\]

and

\[
\boxed{
\text{information flow}.
}
\]

Removing direction can change network conclusions.

%%%%%%%%%%%%%%%%%%%%%%%%%%%%%%%%%%%%%%%%%%%%%%%%%%%%%%%%%%%%
\subsubsection{Ignoring Edge Weights}

A network with one interaction per year and one interaction per second may
have identical unweighted adjacency.

When interaction magnitude matters, weights should be retained.

%%%%%%%%%%%%%%%%%%%%%%%%%%%%%%%%%%%%%%%%%%%%%%%%%%%%%%%%%%%%
\subsubsection{Assuming Degree Equals Importance}

Degree measures local connectivity.

Other forms of importance may depend on:

\[
\boxed{
\text{distance},
}
\]

\[
\boxed{
\text{bridging},
}
\]

or

\[
\boxed{
\text{neighbor importance}.
}
\]

No centrality is universally best.

%%%%%%%%%%%%%%%%%%%%%%%%%%%%%%%%%%%%%%%%%%%%%%%%%%%%%%%%%%%%
\subsubsection{Comparing Centralities Without Considering Definitions}

Degree, closeness, betweenness, eigenvector centrality, Katz centrality, and
PageRank answer different structural questions.

Rankings need not agree.

%%%%%%%%%%%%%%%%%%%%%%%%%%%%%%%%%%%%%%%%%%%%%%%%%%%%%%%%%%%%
\subsubsection{Ignoring Disconnected Components in Closeness}

Ordinary closeness formulas can fail or become ambiguous when distances are
infinite.

Appropriate harmonic or component-based variants may be needed.

%%%%%%%%%%%%%%%%%%%%%%%%%%%%%%%%%%%%%%%%%%%%%%%%%%%%%%%%%%%%
\subsubsection{Assuming High Betweenness Means High Degree}

Bridge nodes can have low degree but high betweenness.

Local connectivity and global brokerage are different concepts.

%%%%%%%%%%%%%%%%%%%%%%%%%%%%%%%%%%%%%%%%%%%%%%%%%%%%%%%%%%%%
\subsubsection{Assuming All Heavy-Tailed Data Follow a Power Law}

A broad degree distribution does not by itself establish

\[
P(k)\propto k^{-\gamma}.
\]

Competing distributions and statistical goodness-of-fit should be examined.

%%%%%%%%%%%%%%%%%%%%%%%%%%%%%%%%%%%%%%%%%%%%%%%%%%%%%%%%%%%%
\subsubsection{Calling Every Network Scale-Free}

The term should be used only when the relevant scaling behavior is actually
supported.

Visual appearance on a log-log plot is not sufficient evidence.

%%%%%%%%%%%%%%%%%%%%%%%%%%%%%%%%%%%%%%%%%%%%%%%%%%%%%%%%%%%%
\subsubsection{Assuming Community Detection Has One Correct Answer}

Communities can depend on:

\[
\boxed{
\text{resolution},
}
\]

\[
\boxed{
\text{algorithm},
}
\]

\[
\boxed{
\text{edge definition},
}
\]

and

\[
\boxed{
\text{null model}.
}
\]

Community structure is often scale-dependent.

%%%%%%%%%%%%%%%%%%%%%%%%%%%%%%%%%%%%%%%%%%%%%%%%%%%%%%%%%%%%
\subsubsection{Treating Modularity as an Absolute Measure}

Modularity is defined relative to a null model.

Its magnitude should be interpreted in context rather than as a universal
quality score.

%%%%%%%%%%%%%%%%%%%%%%%%%%%%%%%%%%%%%%%%%%%%%%%%%%%%%%%%%%%%
\subsubsection{Ignoring Temporal Order}

Aggregating a temporal network can create paths that never existed in a valid
time order.

Temporal reachability can differ substantially from static reachability.

%%%%%%%%%%%%%%%%%%%%%%%%%%%%%%%%%%%%%%%%%%%%%%%%%%%%%%%%%%%%
\subsubsection{Ignoring Sampling Bias}

Observed networks may omit isolated nodes, weak ties, or entire communities.

Computed network statistics may therefore be biased.

%%%%%%%%%%%%%%%%%%%%%%%%%%%%%%%%%%%%%%%%%%%%%%%%%%%%%%%%%%%%
\subsubsection{Treating Missing Edges as Confirmed Nonedges}

Failure to observe an interaction does not prove the interaction is absent.

Measurement processes matter.

%%%%%%%%%%%%%%%%%%%%%%%%%%%%%%%%%%%%%%%%%%%%%%%%%%%%%%%%%%%%
\subsubsection{Confusing Correlation With Network Interaction}

Correlation can reflect indirect pathways or common drivers.

An inferred correlation network does not automatically identify direct
causal edges.

%%%%%%%%%%%%%%%%%%%%%%%%%%%%%%%%%%%%%%%%%%%%%%%%%%%%%%%%%%%%
\subsubsection{Assuming Static Network Structure Determines Dynamics Alone}

Spreading and synchronization also depend on:

\[
\boxed{
\text{rate parameters},
}
\]

\[
\boxed{
\text{initial conditions},
}
\]

\[
\boxed{
\text{nonlinear dynamics},
}
\]

and

\[
\boxed{
\text{edge timing}.
}
\]

Topology is only one part of the model.

%%%%%%%%%%%%%%%%%%%%%%%%%%%%%%%%%%%%%%%%%%%%%%%%%%%%%%%%%%%%
\subsubsection{Confusing Random Walks With Uniform Node Sampling}

A simple random walk on an undirected graph visits nodes in proportion to
degree at stationarity.

It therefore oversamples high-degree nodes relative to uniform sampling.

%%%%%%%%%%%%%%%%%%%%%%%%%%%%%%%%%%%%%%%%%%%%%%%%%%%%%%%%%%%%
\subsubsection{Ignoring Network Direction in Random Walks}

Directed networks require care because stationary distributions may not have
the simple degree-proportional form.

Dangling nodes and non-strong connectivity can also matter.

%%%%%%%%%%%%%%%%%%%%%%%%%%%%%%%%%%%%%%%%%%%%%%%%%%%%%%%%%%%%
\subsubsection{Assuming Robustness to Random Failure Implies Robustness to Targeted Attack}

A network can tolerate random removals while being highly sensitive to
removal of hubs or critical bridges.

The failure mechanism matters.

%%%%%%%%%%%%%%%%%%%%%%%%%%%%%%%%%%%%%%%%%%%%%%%%%%%%%%%%%%%%
\subsubsection{Ignoring Cascading Effects}

Removing one node may change loads or dynamics elsewhere.

Independent node-removal models can miss cascading failures.

%%%%%%%%%%%%%%%%%%%%%%%%%%%%%%%%%%%%%%%%%%%%%%%%%%%%%%%%%%%%
\subsubsection{Confusing Structural and Functional Connectivity}

An edge may exist structurally but carry negligible functional flow.

Conversely, functional dependence may occur through indirect pathways.

%%%%%%%%%%%%%%%%%%%%%%%%%%%%%%%%%%%%%%%%%%%%%%%%%%%%%%%%%%%%
\subsubsection{Assuming More Edges Always Improve a Network}

Additional edges may improve redundancy but can also:

\[
\boxed{
\text{spread contagion faster},
}
\]

\[
\boxed{
\text{increase systemic coupling},
}
\]

or

\[
\boxed{
\text{create new cascade pathways}.
}
\]

Connectivity has both benefits and risks.

%%%%%%%%%%%%%%%%%%%%%%%%%%%%%%%%%%%%%%%%%%%%%%%%%%%%%%%%%%%%
\subsection{Applications}
\label{subsec:network-science-applications}
%%%%%%%%%%%%%%%%%%%%%%%%%%%%%%%%%%%%%%%%%%%%%%%%%%%%%%%%%%%%

\begin{applicationbox}

\textbf{Social Systems.}
Network methods study friendship, collaboration, communication, influence,
communities, and information diffusion.

\medskip

\textbf{Transportation.}
Road, airline, shipping, and transit networks are analyzed using paths,
centrality, flow, robustness, and optimization.

\medskip

\textbf{Communication Systems.}
Routing, reliability, congestion, and information transfer depend on network
topology.

\medskip

\textbf{Biology.}
Gene regulation, protein interactions, metabolism, neural systems, and
ecological interactions naturally form networks.

\medskip

\textbf{Epidemiology.}
Contact networks determine who can transmit infection to whom and can alter
epidemic thresholds and intervention effects.

\medskip

\textbf{Finance.}
Interbank loans, payment systems, asset overlap, and counterparty exposures
create networks of systemic risk.

\medskip

\textbf{Infrastructure.}
Power, water, transportation, and communication systems can experience
network failures and cascades.

\medskip

\textbf{Control and Robotics.}
Network Laplacians describe consensus, formation control, synchronization,
and distributed estimation.

\end{applicationbox}

%%%%%%%%%%%%%%%%%%%%%%%%%%%%%%%%%%%%%%%%%%%%%%%%%%%%%%%%%%%%
\subsection{Exercises}
\label{subsec:network-science-exercises}
%%%%%%%%%%%%%%%%%%%%%%%%%%%%%%%%%%%%%%%%%%%%%%%%%%%%%%%%%%%%

\begin{exercise}[1]
Define a network mathematically as a graph.
\end{exercise}

\begin{exercise}[1]
Distinguish directed and undirected networks.
\end{exercise}

\begin{exercise}[1]
Define a weighted network.
\end{exercise}

\begin{exercise}[1]
Define the adjacency matrix.
\end{exercise}

\begin{exercise}[1]
Define degree.
\end{exercise}

\begin{exercise}[1]
Define in-degree and out-degree.
\end{exercise}

\begin{exercise}[1]
Define graph distance.
\end{exercise}

\begin{exercise}[1]
Define a connected component.
\end{exercise}

\begin{exercise}[1]
Define a clustering coefficient.
\end{exercise}

\begin{exercise}[1]
Define centrality conceptually.
\end{exercise}

\begin{exercise}[2]
Construct the adjacency matrix for a four-node cycle.
\end{exercise}

\begin{exercise}[2]
Compute all node degrees in a four-node path.
\end{exercise}

\begin{exercise}[2]
Verify the handshaking relation for a given undirected graph.
\end{exercise}

\begin{exercise}[2]
Compute average degree from \(|V|\) and \(|E|\).
\end{exercise}

\begin{exercise}[2]
Find all connected components of a given graph.
\end{exercise}

\begin{exercise}[2]
Compute shortest-path distances in a small network.
\end{exercise}

\begin{exercise}[2]
Compute the diameter of a path graph with \(n\) vertices.
\end{exercise}

\begin{exercise}[2]
Compute the local clustering coefficient of a specified node.
\end{exercise}

\begin{exercise}[3]
Show that

\[
\frac16\operatorname{tr}(A^3)
\]

counts triangles in a simple undirected graph.
\end{exercise}

\begin{exercise}[2]
Compute degree centrality for each node in a small graph.
\end{exercise}

\begin{exercise}[2]
Compute closeness centrality in a connected four-node graph.
\end{exercise}

\begin{exercise}[3]
Compute betweenness centrality for a path graph.
\end{exercise}

\begin{exercise}[3]
Explain why the center of a path has high betweenness.
\end{exercise}

\begin{exercise}[2]
Explain eigenvector centrality conceptually.
\end{exercise}

\begin{exercise}[3]
Compute the dominant eigenvector of a small adjacency matrix.
\end{exercise}

\begin{exercise}[2]
Construct the degree matrix \(D\).
\end{exercise}

\begin{exercise}[2]
Construct the graph Laplacian

\[
L=D-A.
\]
\end{exercise}

\begin{exercise}[3]
Verify

\[
L\mathbf{1}=0.
\]
\end{exercise}

\begin{exercise}[3]
Verify the Laplacian quadratic-form identity for a small graph.
\end{exercise}

\begin{exercise}[3]
Compute Laplacian eigenvalues of a three-node path.
\end{exercise}

\begin{exercise}[3]
Identify algebraic connectivity.
\end{exercise}

\begin{exercise}[3]
Explain how a Fiedler vector can partition a graph.
\end{exercise}

\begin{exercise}[2]
Define the Erdős--Rényi \(G(n,p)\) model.
\end{exercise}

\begin{exercise}[2]
Compute expected number of edges in \(G(n,p)\).
\end{exercise}

\begin{exercise}[2]
Compute expected degree in \(G(n,p)\).
\end{exercise}

\begin{exercise}[3]
Derive the binomial degree distribution of \(G(n,p)\).
\end{exercise}

\begin{exercise}[3]
Explain the giant-component transition conceptually.
\end{exercise}

\begin{exercise}[2]
Explain the small-world property.
\end{exercise}

\begin{exercise}[2]
Explain preferential attachment.
\end{exercise}

\begin{exercise}[3]
Explain why preferential attachment can produce hubs.
\end{exercise}

\begin{exercise}[3]
Explain why a heavy-tailed degree distribution does not automatically prove
a power law.
\end{exercise}

\begin{exercise}[2]
Explain the configuration model.
\end{exercise}

\begin{exercise}[2]
Define assortative mixing.
\end{exercise}

\begin{exercise}[2]
Define homophily.
\end{exercise}

\begin{exercise}[3]
Explain why homophily and social influence can be difficult to distinguish
from observational network data.
\end{exercise}

\begin{exercise}[2]
Define network community conceptually.
\end{exercise}

\begin{exercise}[2]
Explain modularity conceptually.
\end{exercise}

\begin{exercise}[3]
Compute modularity for a small proposed partition.
\end{exercise}

\begin{exercise}[3]
Explain stochastic block models.
\end{exercise}

\begin{exercise}[2]
Define a bipartite network.
\end{exercise}

\begin{exercise}[3]
Construct a one-mode projection of a small bipartite graph.
\end{exercise}

\begin{exercise}[3]
Explain information lost during bipartite projection.
\end{exercise}

\begin{exercise}[2]
Define a multilayer network.
\end{exercise}

\begin{exercise}[2]
Define a temporal network.
\end{exercise}

\begin{exercise}[3]
Explain a time-respecting path.
\end{exercise}

\begin{exercise}[2]
Define site percolation.
\end{exercise}

\begin{exercise}[2]
Define bond percolation.
\end{exercise}

\begin{exercise}[3]
Explain how percolation models network robustness.
\end{exercise}

\begin{exercise}[3]
Compare random failures with targeted node removal.
\end{exercise}

\begin{exercise}[3]
Explain cascading failure conceptually.
\end{exercise}

\begin{exercise}[2]
Construct a random-walk transition matrix from an adjacency matrix.
\end{exercise}

\begin{exercise}[3]
Find the stationary distribution of a random walk on a small undirected
graph.
\end{exercise}

\begin{exercise}[3]
Show that

\[
\pi_i
=
\frac{k_i}{2|E|}
\]

is stationary for a connected undirected graph.
\end{exercise}

\begin{exercise}[2]
Write the network diffusion equation

\[
\dot{\mathbf{x}}
=
-\kappa L\mathbf{x}.
\]
\end{exercise}

\begin{exercise}[3]
Solve network diffusion using a matrix exponential.
\end{exercise}

\begin{exercise}[3]
Show that total state

\[
\mathbf{1}^T\mathbf{x}
\]

is conserved under undirected Laplacian diffusion.
\end{exercise}

\begin{exercise}[2]
Write the consensus equation.
\end{exercise}

\begin{exercise}[3]
Explain the role of \(\lambda_2\) in consensus convergence.
\end{exercise}

\begin{exercise}[2]
Write the Kuramoto network model.
\end{exercise}

\begin{exercise}[3]
Explain synchronization conceptually.
\end{exercise}

\begin{exercise}[3]
Explain why network topology can affect synchronization.
\end{exercise}

\begin{exercise}[2]
Explain epidemic spreading on a contact network.
\end{exercise}

\begin{exercise}[3]
Explain why adjacency eigenvalues can influence network epidemic thresholds.
\end{exercise}

\begin{exercise}[2]
Define a threshold contagion model.
\end{exercise}

\begin{exercise}[3]
Construct a small example of a cascade triggered by one initially active
node.
\end{exercise}

\begin{exercise}[2]
Define a flow network.
\end{exercise}

\begin{exercise}[2]
State flow-conservation constraints.
\end{exercise}

\begin{exercise}[2]
Define a network cut.
\end{exercise}

\begin{exercise}[2]
State the max-flow min-cut theorem.
\end{exercise}

\begin{exercise}[3]
Solve a small maximum-flow problem.
\end{exercise}

\begin{exercise}[3]
Find a corresponding minimum cut.
\end{exercise}

\begin{exercise}[2]
Formulate a shortest-path problem on a weighted network.
\end{exercise}

\begin{exercise}[3]
Explain network resilience.
\end{exercise}

\begin{exercise}[3]
Give an example of structural connectivity failing to represent functional
resilience.
\end{exercise}

\begin{exercise}[2]
Give an example of a biological network.
\end{exercise}

\begin{exercise}[2]
Give an example of a financial network.
\end{exercise}

\begin{exercise}[2]
Give an example of an interdependent infrastructure network.
\end{exercise}

\begin{exercise}[2]
Define a network motif.
\end{exercise}

\begin{exercise}[3]
Explain why motif significance requires a null model.
\end{exercise}

\begin{exercise}[2]
Define network inference.
\end{exercise}

\begin{exercise}[2]
Define link prediction.
\end{exercise}

\begin{exercise}[2]
Compute a common-neighbor link score.
\end{exercise}

\begin{exercise}[3]
Explain how random-walk network sampling biases observed degree.
\end{exercise}

\begin{exercise}[3]
Explain why missing edges create uncertainty.
\end{exercise}

\begin{exercise}[3]
Explain why correlation networks should not automatically be interpreted
causally.
\end{exercise}

\begin{exercise}[2]
Compute entropy of a simple degree distribution.
\end{exercise}

\begin{exercise}[3]
Distinguish dynamics on networks from dynamics of networks.
\end{exercise}

\begin{exercise}[3]
Construct a simple adaptive-network model conceptually.
\end{exercise}

\begin{exercise}[3]
Explain network controllability.
\end{exercise}

\begin{exercise}[3]
Explain how actuator placement changes controllability.
\end{exercise}

\begin{exercise}[3]
Explain network observability and sensor placement.
\end{exercise}

\begin{exercise}[4]
Analyze Laplacian spectra of several graph families.
\end{exercise}

\begin{exercise}[4]
Derive Laplacian eigenvalues of a complete graph.
\end{exercise}

\begin{exercise}[4]
Derive Laplacian eigenvalues of a cycle graph.
\end{exercise}

\begin{exercise}[4]
Implement spectral clustering on a weighted graph.
\end{exercise}

\begin{exercise}[4]
Simulate an Erdős--Rényi phase transition.
\end{exercise}

\begin{exercise}[4]
Compare random, small-world, and preferential-attachment networks using
degree, clustering, and path-length statistics.
\end{exercise}

\begin{exercise}[4]
Fit and statistically evaluate a candidate power-law degree distribution.
\end{exercise}

\begin{exercise}[4]
Implement a configuration-model network.
\end{exercise}

\begin{exercise}[4]
Study modularity maximization and its resolution limitations.
\end{exercise}

\begin{exercise}[4]
Fit a stochastic block model conceptually or computationally.
\end{exercise}

\begin{exercise}[4]
Study percolation under random and targeted node removal.
\end{exercise}

\begin{exercise}[4]
Simulate a cascading-failure model.
\end{exercise}

\begin{exercise}[4]
Study mixing times of random walks on graphs.
\end{exercise}

\begin{exercise}[4]
Analyze consensus convergence using Laplacian eigenvalues.
\end{exercise}

\begin{exercise}[4]
Simulate synchronization in a Kuramoto network.
\end{exercise}

\begin{exercise}[4]
Compare homogeneous-mixing and network epidemic models.
\end{exercise}

\begin{exercise}[4]
Study epidemic thresholds for networks with different degree distributions.
\end{exercise}

\begin{exercise}[4]
Simulate threshold cascades on clustered networks.
\end{exercise}

\begin{exercise}[4]
Solve a maximum-flow problem using linear programming.
\end{exercise}

\begin{exercise}[4]
Analyze resilience of a transportation network.
\end{exercise}

\begin{exercise}[4]
Construct and analyze a temporal network.
\end{exercise}

\begin{exercise}[4]
Develop a link-prediction model and evaluate it on held-out edges.
\end{exercise}

\begin{exercise}[4]
Study network controllability for several actuator placements.
\end{exercise}

\begin{exercise}[4]
Study sensor placement and observability on a networked dynamical system.
\end{exercise}

%%%%%%%%%%%%%%%%%%%%%%%%%%%%%%%%%%%%%%%%%%%%%%%%%%%%%%%%%%%%
\subsection{Conceptual Summary}
\label{subsec:network-science-summary}
%%%%%%%%%%%%%%%%%%%%%%%%%%%%%%%%%%%%%%%%%%%%%%%%%%%%%%%%%%%%

Network science studies systems through the mathematical structure of their
connections.

A network is represented by

\[
\boxed{
G=(V,E).
}
\]

The adjacency matrix

\[
\boxed{
A
}
\]

encodes direct relationships.

Local connectivity is measured by degree:

\[
\boxed{
k_i
=
\sum_jA_{ij}.
}
\]

Global geometry is described by paths and distances.

Clustering describes local closure.

Centrality measures different forms of structural importance.

The graph Laplacian

\[
\boxed{
L=D-A
}
\]

connects network structure with linear algebra and dynamics.

Its spectrum reveals information about connectivity and diffusion.

Random graph models provide baselines for expected network structure.

Small-world, preferential-attachment, block, temporal, and multilayer models
capture additional forms of organization.

Processes such as

\[
\boxed{
\text{random walks},
}
\]

\[
\boxed{
\text{diffusion},
}
\]

\[
\boxed{
\text{consensus},
}
\]

\[
\boxed{
\text{synchronization},
}
\]

\[
\boxed{
\text{contagion},
}
\]

and

\[
\boxed{
\text{flow}
}
\]

evolve on networks.

Thus network science combines

\[
\boxed{
\text{graph theory}
+
\text{linear algebra}
+
\text{probability}
+
\text{dynamical systems}
+
\text{statistics}
+
\text{optimization}.
}
\]

The central applied question is not merely which nodes and edges exist, but

\[
\boxed{
\text{how network structure shapes system behavior}.
}
\]

%%%%%%%%%%%%%%%%%%%%%%%%%%%%%%%%%%%%%%%%%%%%%%%%%%%%%%%%%%%%
\subsection{Section Connections}
\label{subsec:network-science-connections}
%%%%%%%%%%%%%%%%%%%%%%%%%%%%%%%%%%%%%%%%%%%%%%%%%%%%%%%%%%%%

\chapterconnection{
Network Science expands the graph-theoretic foundations of Chapter~6 into a
general framework for applied systems whose interactions occur through
networks. Probability provides random graph and random-walk models, linear
algebra provides adjacency and Laplacian spectral methods, dynamical systems
describe diffusion and synchronization, Control Theory provides consensus,
controllability, and observability, and optimization provides routing,
maximum-flow, and network-design methods. Network models also connect
directly to epidemiology, ecology, infrastructure, communication, social
systems, and finance. The next section, Mathematical Finance, develops
mathematical models of asset prices, returns, portfolios, interest rates,
derivatives, arbitrage, risk, stochastic processes, and financial decision
making.
}

%%%%%%%%%%%%%%%%%%%%%%%%%%%%%%%%%%%%%%%%%%%%%%%%%%%%%%%%%%%%
% SECTION SOURCE TRACKING
%%%%%%%%%%%%%%%%%%%%%%%%%%%%%%%%%%%%%%%%%%%%%%%%%%%%%%%%%%%%

%------------------------------------------------------------
% 13.7 Network Science
%------------------------------------------------------------
%
% Source basis:
%   - Standard graph theory and network science.
%   - Standard random graph, spectral graph, diffusion,
%     centrality, community, robustness, contagion,
%     flow, and network-dynamics methods.
%
% Used for:
%   - networks and graphs;
%   - directed and undirected networks;
%   - weighted networks;
%   - adjacency matrices;
%   - degree, in-degree, out-degree, and strength;
%   - degree distributions;
%   - sparse networks;
%   - paths and distances;
%   - connectivity;
%   - diameter and average path length;
%   - triangles and clustering;
%   - centrality;
%   - degree centrality;
%   - closeness;
%   - betweenness;
%   - eigenvector centrality;
%   - PageRank;
%   - Katz centrality;
%   - graph Laplacians;
%   - normalized Laplacians;
%   - Laplacian spectra;
%   - algebraic connectivity;
%   - Fiedler vectors;
%   - spectral clustering;
%   - Erdős--Rényi random graphs;
%   - expected degree;
%   - random graph degree distributions;
%   - giant components;
%   - small-world networks;
%   - Watts--Strogatz models;
%   - preferential attachment;
%   - scale-free networks;
%   - configuration models;
%   - assortativity;
%   - homophily;
%   - communities;
%   - modularity;
%   - stochastic block models;
%   - hierarchical networks;
%   - bipartite networks;
%   - multilayer and multiplex networks;
%   - temporal networks;
%   - robustness;
%   - site and bond percolation;
%   - cascading failure;
%   - random walks;
%   - stationary distributions;
%   - network diffusion;
%   - consensus;
%   - synchronization;
%   - Kuramoto networks;
%   - network epidemics;
%   - spectral epidemic thresholds;
%   - threshold contagion;
%   - flow networks;
%   - maximum flow;
%   - minimum cut;
%   - shortest paths;
%   - resilience;
%   - social networks;
%   - biological and ecological networks;
%   - financial and infrastructure networks;
%   - interdependent networks;
%   - motifs;
%   - null models;
%   - network inference;
%   - link prediction;
%   - sampling;
%   - missing edges;
%   - correlation networks;
%   - network entropy;
%   - dynamic and adaptive networks;
%   - network controllability;
%   - network observability;
%   - network optimization;
%   - applications and exercises.
%
% Notes:
%   - Mathematical Finance is developed next in Section 13.8.
%   - Mathematical Economics follows in Section 13.9.
%
%%%%%%%%%%%%%%%%%%%%%%%%%%%%%%%%%%%%%%%%%%%%%%%%%%%%%%%%%%%%